% --- CHAPTER 9: SELF-EVALUATION ---
\section{Self-Evaluation}
\label{sec:self_evaluation}

This chapter contains the individual self-evaluation statements for each group member. In accordance with the academic guidelines, each section outlines individual responsibilities, objective role assessments, and a critical analysis of the product's engineering strengths and structural weaknesses from each developer's perspective.

\subsection{Sajmir Buzi}
\label{subsec:eval_sajmir}
Working on the design and development of AdriaBOX has been a deeply educational experience.
Developing a system from scratch allowed me to face the real-world challenges of distributed
systems, forcing me to turn the theoretical concepts learned in class into practical solutions.
In particular, it was both challenging and rewarding to manage real memory constraints—such
as the 4KB page limit for streaming—and to guarantee fault tolerance without compromising
data consistency when nodes fail.

A fundamental strength of this project was the workflow I established with my colleague.
We carefully planned the architecture right from the start, clearly defining how the different
components would communicate with each other. This approach allowed us to write the code
smoothly, preventing any major roadblocks or compatibility issues during integration.

We are so satisfied with the final result that we are considering looking beyond the academic
scope: our idea for the future is to further engineer the product for potential commercialization.
The first step in this direction will be developing a proper graphical user interface (GUI) to
make AdriaBOX accessible and scalable even for non-technical users, turning a solid backend
into a complete cloud service.

\subsection{Massimiliano Pupillo}
\label{subsec:eval_massimiliano}
The experience of designing and deploying the AdriaBOX distributed architecture has represented a cornerstone in my academic and professional development. The most significant element of this project lies not only in the technical implementation of the distributed storage logic, but in the specific collaborative methodology established with my colleague. From the very inception of the infrastructure, our workflow was never partitioned in a dichotomous or fragmented manner through a rigid separation of modules; instead, we worked strictly side-by-side across every phase of the project, co-engineering everything from low-level network interactions to high-level storage persistence mappings.

Given our fundamentally different academic backgrounds—with Sajmir rooted in pure computer science and my trajectory anchored in electronic engineering—each of us approached architectural roadblocks from entirely distinct engineering perspectives. What truly formed and enriched my technical paradigm was this continuous exposure to and sharing of divergent problem-solving strategies. My perspective, naturally focused on physical hardware constraints, low-level virtual memory mapping configurations, and structural resource boundary optimization, integrated seamlessly with Sajmir's high-level software abstraction patterns and algorithmic reasoning. This mutual cross-contamination allowed us to discover innovative technical paths and analyze systemic efficiency at a level of granularity that neither of us could have reached in isolation.

The professional and personal synergy resulting from this collaboration has been exceptional. We are so satisfied with the computational resilience and core stability of the AdriaBOX backend that we are actively exploring market avenues for potential commercialization, with the shared goal of building a technological venture and engineering future systems together.